% ==============================================================================
% Fichier : 04_referentiels.tex
% Description : Normes et Référentiels Usuels en Audit des SI
% ==============================================================================

\chapter{Normes et Référentiels Usuels en Audit des Systèmes d'Information}\label{chap:04_referentiels}

L'auditeur a besoin de "briques" techniques et organisationnelles pour évaluer la maturité d'un SI. Les référentiels fournissent les bonnes pratiques reconnues internationalement. Ce chapitre présente les trois piliers : COSO, COBIT et la série ISO 27000.

\section{COSO : Le Modèle de Contrôle Interne}
% ------------------------------------------------------------------------------

\subsection{Présentation}

\begin{boxdef}
\textbf{COSO (Committee of Sponsoring Organizations) :} Cadre de référence conceptuel du contrôle interne, publié en 1992 (révisé en 2013). C'est le modèle standard pour la gouvernance d'entreprise, souvent exigé par les lois sur la gouvernance (ex: Sarbanes-Oxley aux USA).
\end{boxdef}

Bien que généraliste, COSO est le socle sur lequel s'appuient les autres référentiels informatiques.

\subsection{Les Trois Objectifs}

Le contrôle interne selon COSO vise à donner une assurance raisonnable sur la réalisation de trois objectifs :
\begin{enumerate}
    \item \textbf{Réalisation et Optimisation des opérations} (Performance).
    \item \textbf{Fiabilité des informations financières} (Comptabilité).
    \item \textbf{Conformité aux lois et règlements} (Juridique).
\end{enumerate}

\subsection{Les Cinq Composantes (Le Cube COSO)}

Pour atteindre ces objectifs, l'organisation doit mettre en place 5 composantes liées :

\begin{table}[h!]
\centering
\footnotesize
\begin{tabular}{p{3.5cm}p{7cm}}
    \toprule
    \textbf{Composante} & \textbf{Description} \\
    \midrule
    \textbf{Environnement} & Culture, éthique, compétence du personnel. \\
    \textbf{Évaluation Risques} & Identification des risques empêchant atteindre objectifs. \\
    \textbf{Activités Contrôle} & Politiques et procédures (Séparation tâches, autorisations). \\
    \textbf{Info et Communication} & Circulation d'info nécessaire au pilotage. \\
    \textbf{Pilotage (Monitoring)} & Surveillance continue du système de contrôle. \\
    \bottomrule
\end{tabular}
\caption{Les 5 composantes du référentiel COSO}
\end{table}

\begin{boxlizbiz}
\textbf{Application Lizbiz's :}
L'audit du système de gestion des commandes chez Lizbiz's vérifiera si les 5 composantes COSO sont présentes : existe-t-il une procédure (Activité de contrôle) ? Les risques de fraude sont-ils évalués (Évaluation des risques) ? Les managers supervisent-ils les anomalies (Pilotage) ?
\end{boxlizbiz}

% ------------------------------------------------------------------------------
\section{COBIT : La Gouvernance des SI}
% ------------------------------------------------------------------------------

\subsection{Présentation}

\begin{boxdef}
\textbf{COBIT (Control Objectives for Information and Related Technologies) :} Référentiel de gouvernance et de management des SI développé par l'ISACA. Il fait le pont entre les métiers et l'IT.
\end{boxdef}

C'est le référentiel de choix pour l'auditeur SI car il propose une liste d'objectifs de contrôle précis pour chaque processus informatique.

\subsection{Objectifs et Apports}

COBIT répond à la question : "Comment aligner l'IT sur la stratégie de l'entreprise ?".
Il structure le SI en domaines et processus, offrant à l'auditeur une grille de lecture universelle.

\subsection{Structure de COBIT (Version 5 / 2019)}

COBIT divise la gouvernance et le management en 5 domaines principaux :

\begin{table}[h!]
\centering
\footnotesize
\setlength{\tabcolsep}{3pt}
\begin{tabular}{p{3cm}p{8cm}}
    \toprule
    \textbf{Domaine} & \textbf{Contenu (Exemples)} \\
    \midrule
    \textbf{Evaluer, Diriger et Surveiller (EDM)} & Gouvernance stratégique. \\
    \midrule
    \textbf{Planifier et Organiser (PO)} & Stratégie, Architecture, Budget. \\
    \midrule
    \textbf{Acquérir et Implémenter (AI)} & Projets, Développement, Intégration. \\
    \midrule
    \textbf{Délivrer et Servir (DS)} & Exploitation, Support, Sécurité. \\
    \midrule
    \textbf{Surveiller, Évaluer et Analyser (MEA)} & Performance, Conformité. \\
    \bottomrule
\end{tabular}
\caption{Les domaines de COBIT}
\end{table}

\subsection{Utilisation par l'Auditeur}

Pour chaque processus (ex: DS5 "Gérer la sécurité"), COBIT fournit :
\begin{itemize}
    \item Un objectif de contrôle détaillé.
    \item Des indicateurs de performance (KPI).
    \item Un modèle de maturité (de 0 : Inexistant à 5 : Optimisé).
\end{itemize}

\subsection{La Cascade d'Objectifs COBIT}
L'un des points forts de COBIT est sa capacité à lier la stratégie métier aux objectifs techniques. 
\begin{enumerate}
    \item \textbf{Besoins des parties prenantes} (Bénéfices, Risques, Ressources).
    \item \textbf{Objectifs d'entreprise} (ex: Portefeuille de produits et services structurés).
    \item \textbf{Objectifs d'alignement} (ex: Qualité des informations financières liées à l'IT).
\end{enumerate}
% ------------------------------------------------------------------------------
\section{La Série ISO 27000 : La Sécurité Normative}
% ------------------------------------------------------------------------------

Alors que COBIT gère l'IT, la série ISO 27000 se concentre spécifiquement sur la \textbf{Sécurité de l'Information}.

\subsection{Vue d'Ensemble}

C'est une famille de normes internationales. Les plus critiques pour l'auditeur sont :

\begin{itemize}
    \item \textbf{ISO 27001 :} Exigences pour un Système de Management de la Sécurité de l'Information (SMSI). C'est la norme \textit{certifiable}.
    \item \textbf{ISO 27002 :} Code de bonne pratique. Liste détaillée des mesures de sécurité (contrôles) à implémenter.
    \item \textbf{ISO 27005 :} Gestion des risques en sécurité de l'information.
    \item \textbf{ISO 27006 :} Exigences pour les organismes de certification.
\end{itemize}

\subsection{Importance pour l'Audit}

\begin{enumerate}
    \item \textbf{Critères d'audit :} ISO 27001 fournit les exigences ("Il faut une politique de sécurité"). L'auditeur vérifie si c'est fait.
    \item \textbf{Mesures techniques :} ISO 27002 aide l'auditeur à proposer des recommandations techniques précises (ex: "Comment configurer l'authentification ?").
\end{enumerate}

\begin{boxlizbiz}
\textbf{Projet Lizbiz's :}
Suite à l'audit, la direction de Lizbiz's souhaite rassurer ses clients. Elle décide de se faire certifier \textbf{ISO 27001}. L'auditeur utilisera l'Annexe A de cette norme (qui liste les contrôles obligatoires) comme grille d'évaluation pour le futur audit de certification.
\end{boxlizbiz}

% ------------------------------------------------------------------------------
\section{Synthèse : Comment Choisir ?}
% ------------------------------------------------------------------------------

L'auditeur combine souvent ces référentiels :

\begin{table}[h!]
\centering
\footnotesize
\setlength{\tabcolsep}{3pt}
\begin{tabular}{p{2.8cm}p{3.8cm}p{4.2cm}}
    \toprule
    \textbf{Référentiel} & \textbf{Angle} & \textbf{Question clé} \\
    \midrule
    \textbf{COSO} & Gouvernance Globale & Contrôles bien conçus ? \\
    \textbf{COBIT} & Processus IT & IT bien gérée ? \\
    \textbf{ISO 27001} & Sécurité Info & Sécurisé et certifiable ? \\
    \bottomrule
\end{tabular}
\caption{Comparatif des référentiels}
\end{table}

\begin{table}[h!]
\centering
\footnotesize
\setlength{\tabcolsep}{2pt}
\renewcommand{\arraystretch}{1.2}
\begin{tabular}{|p{1.7cm}|p{1.9cm}|p{2cm}|p{2.2cm}|}
    \hline
    \rowcolor{blue!10} \textbf{Référentiel} & \textbf{Cible} & \textbf{Focus Audit} & \textbf{Livrable} \\ \hline
    \textbf{COSO} & DG / Audit & Contrôle Interne & Rapport Sarbanes-Oxley \\ \hline
    \textbf{COBIT} & DSI / Métiers & Gouvernance IT & Matrice maturité \\ \hline
    \textbf{ISO 27001} & RSSI / Sécu & Sécurité Info & Cert. SMSI \\ \hline
\end{tabular}
\caption{Tableau comparatif des référentiels d'audit}
\end{table}
